\chapter{Saturation et vibrance : ce que l'on cherche à reproduire}

\section{Deux curseurs voisins, deux intentions}
Les deux curseurs modifient l'intensité des couleurs, sans changer la
luminosité générale de l'image. Leur intention diffère pourtant.

\begin{itemize}
\item La \textbf{saturation} agit uniformément : chaque couleur voit son
intensité multipliée par le même facteur. À $-100$, l'image devient grise ; à
$+100$, toutes les couleurs, y compris celles qui sont déjà vives, sont
renforcées autant.
\item La \textbf{vibrance} est sélective. Elle renforce surtout les couleurs
ternes et épargne celles qui sont déjà intenses. Les tons chair sont réputés
protégés. Le photographe s'en sert pour réveiller une image sans faire
« baver » les couleurs vives.
\end{itemize}

La description de la vibrance reste qualitative : les éditeurs ne publient pas
leur formule. Pour reproduire Lightroom, il faut donc le \emph{mesurer}.

\section{Pourquoi un modèle mathématique simple ne suffit pas}
Un premier modèle de vibrance, utilisé par \IPR{} jusqu'en octobre 2026,
pondérait le gain par la saturation relative du pixel,
$g_v=1+(V/100)(1-c)$ avec $c=(\max-\min)/\max$. Un second prototype (A1/A2)
travaillait en Oklab avec un poids $(1-q)^2$ et une protection des tons chair.
Mesurés sur 49 scènes, ces modèles restaient loin de Lightroom :

\begin{itemize}
\item à $+25$ et $+50$, ils renforçaient environ deux fois trop les couleurs
peu saturées ;
\item à $+100$, ils renforçaient trop les couleurs saturées ;
\item à $-100$, ils retiraient trop d'intensité aux couleurs saturées ;
\item ils ne modifiaient jamais la clarté, alors que Lightroom assombrit
légèrement en désaturant, ni la teinte, alors que Lightroom la fait tourner.
\end{itemize}

Le dernier point est le plus parlant : aucun choix de constantes n'aurait
corrigé un modèle qui ne possède pas le degré de liberté nécessaire.

\resultat{On ne devine pas le comportement d'un logiciel de référence : on le
mesure, on identifie les variables qui l'expliquent, puis on choisit la
représentation du modèle.}

\section{La démarche suivie}
\begin{enumerate}
\item \textbf{Constituer des paires} : pour chaque photographie, un export
Lightroom neutre et un export avec le seul curseur modifié.
\item \textbf{Comparer pixel par pixel} chaque export à son neutre : gain
d'intensité, changement de clarté, rotation de teinte.
\item \textbf{Chercher les variables explicatives} : de quoi dépend la réponse
d'un pixel ? De sa saturation, de sa clarté, de sa teinte, de l'image ?
\item \textbf{Construire un modèle} à partir de ces variables.
\item \textbf{Valider} sur des images exclues de la construction, puis sur le
rendu réel d'\IPR{}.
\end{enumerate}

Ces cinq étapes structurent les chapitres suivants.
